\section{Zadanie 12}
Niech \(T_1, T_2\) będą zmiennymi o rozkładzie wykładniczym z parametrem 1 i niech \(T=T_1+T_2\). Oblicz \(E(T_1|T)\) oraz \(E(T|T_1)\).

\subsection*{Rozwiązanie}

W ogólnym przypadku, gdy nie wiemy nic o zależności \(T_1\) oraz \(T_2\), możemy jedynie stwierdzić:
\begin{align*}
	E(T_1|T)(t) & = \int_{-\infty}^{\infty} x \cdot f_{T_1| T}(x,t)dx = \int_{-\infty}^{\infty} x \cdot \frac{f_{T_1, T}{(x, t)}}{f_{T}(t)}dx= \\
	            & = \int_{-\infty}^{\infty} x \cdot \frac{f_{T_1, T_2}{(x, t-x)}}{f_{T}(t)}dx
\end{align*}
Z kolei druga warunkowa wartość oczekiwana:
\begin{align*}
	E(T|T_1) = E(T_1+T_2|T_1) = E(T_1|T_1) + E(T_2|T_1) = T_1 + E(T_2|T_1)
	\\
	E(T|T_1)(t) = t + \int_{-\infty}^{\infty}x \cdot f_{T_2 | T_1}(x, t)dx = t + \int_{-\infty}^{\infty}x\cdot \frac{f_{T_1, T_2}(t, x)}{f_{T_1}(t)} dx
\end{align*}

W przypadku, gdy \(T_1\) i \(T_2\) niekoniecznie są niezależne, nie znając \(f_{T_1, T_2}\), nie jesteśmy w stanie obliczyć konkretnych wartości oczekiwanych. Z kolei gdy \(T_1, T_2\) niezależne, to:\\
\textbf{1.} z symetrii która wynika z niezależności i takiego samego parametru rozkładu wykładniczego: \(E(T_1|T) = E(T_2|T) = T/2\)\\
\textbf{2.} z niezależności: \(E(T|T_1) = T_1 + E(T_2|T_1) = T_1 + E(T_2) = T_1 + \frac{1}{\lambda} = T_1 + 1\)

\subsection*{Odpowiedź}
\[
	\boxed{
		\begin{aligned}
			\text{Dla przypadku } T_1, T_2 \text{ niezależne:} \\
			E(T_1 \mid T) & = \frac{T}{2},                     \\
			E(T \mid T_1) & = T_1 + 1.
		\end{aligned}
	}
\]
